\begin{table}[ht]
\centering
\small
\caption{Model selection for cockle occurrence, biomass where present and both together (the two-part model). $\Delta$AIC is the difference in AIC to the best model (lower AIC is better). $\Delta$ELPD is the difference in expected log predictive density to the best model in spatially blocked 10-fold cross-validation (higher ELPD is better), with its standard error in parentheses. The best model in each column is shown in bold. The total AIC and ELPD are the sums of the two parts. Occurrence is fitted to all surveys, with survey (gear) and year as factors; biomass where present to the grab survey only, with year as a factor and previous-season fishing effort (swept-area ratio of the cockle fishery). Both parts include a spatial random field. Connectivity is presence-weighted in-strength; environment comprises depth (linear and quadratic), temperature, oxygen, salinity, and maximum shear stress, plus sediment type in the occurrence part. A dash marks a model that did not converge.}
\label{tab:model-selection}
\begin{tabular}{lrrrrrr}
\toprule
 & \multicolumn{3}{c}{$\Delta$AIC} & \multicolumn{3}{c}{$\Delta$ELPD (SE)} \\
\cmidrule(lr){2-4} \cmidrule(lr){5-7}
Model & Occurrence & Biomass & Total & Occurrence & Biomass & Total \\
\midrule
Space + Environment + Connectivity & \textbf{0.0} & \textbf{0.0} & \textbf{0.0} & \textbf{0.0} & -7.1 (7.2) & \textbf{0.0} \\
Space + Environment & 2.9 & 3.1 & 6.0 & -22.5 (5.8) & -15.5 (11.0) & -30.9 (9.1) \\
Space + Connectivity & 66.4 & 4.3 & 70.7 & -111.1 (13.4) & \textbf{0.0} & -104.0 (15.0) \\
Space & 65.5 & 10.9 & 76.3 & -115.1 (13.8) & -9.5 (8.1) & -117.5 (16.2) \\
\bottomrule
\end{tabular}
\end{table}
